\subsection{Balancing Intuition and Deliberation in AI Decision-Making}
\label{sec:5.3.2}
If the fast and slow labels do not map cleanly onto untrustworthy and trustworthy, as the previous section argues, an AI system's moral judgment cannot be built on the assumption that whichever answer survives deliberation is thereby the better one. A more defensible balance treats the two processes as different lenses on the same situation, with neither standing as producer to the other's editor.

The fast, pattern-matched read is not noise to be filtered out. It often tracks something a slower calculation would miss or discount — the difference between pushing someone into a trolley's path with one's own hands and throwing a switch at a distance is a difference in what the agent is doing, and the arithmetic being identical is what makes that easy to miss. Some intuitive responses run the other way. The psychologist Paul Slovic's research on aid and disaster response found that people's willingness to help falls as the number of people needing help rises, and that a single named, photographed victim draws more donations than a statistic describing thousands of anonymous ones \autocite{slovic2007mass}. That gap between the vivid and the abstract is the kind of intuitive signal deliberation exists to catch — not because intuition is wrong in general, but because caring less as the stakes grow larger is a well-documented, reliably backward pattern.

The design implication is a system that asks two questions of its own judgment, rather than one process checking the other's work. First, what did the fast read notice that a from-scratch calculation would not have gone looking for — personal agency, directness of harm, a vulnerable party easy to lose in an aggregate? Second, which of deliberation's own known failure modes — scope insensitivity, in-group bias, the identifiable-victim effect — might be distorting the fast read here? Disagreement between the two is a signal to gather more information or escalate the decision, and nothing about it instructs the system to defer to whichever process ran slower.

An autonomous vehicle facing an unavoidable-harm scenario poses the same question where it costs most. A system hard-coded to take whichever option a slower calculation ranks highest has already made the footbridge case's decision by default — the coolest, least-flinching answer wins because it ran slower. The more defensible design specifies in advance which features of a scenario the system should weigh — proximity of harm, the number and vulnerability of the parties involved, whether the vehicle's own action or another actor's created the danger — and requires a fast risk assessment and a check against that specification to agree, or to flag where they do not. That design puts the values into the system ahead of time; the alternative trusts whichever process runs fastest under pressure to supply them. Training a system to run that check is itself a learning problem, exposing it to cases where the two reads disagree and inferring from how people resolve them which situations call for weighting one signal over the other.
